\pg{358} \emph{necessario}, \emph{infinito}, \emph{simplicissimo}, \emph{é immutavel} \footnote{O ente. \emph{immutavel} --- é \emph{necessario}, porque, se fosse \emph{contingente}, receberia de outro ente a. existencia, e. pelo facto de passar da potencia para o acto, já não seria \emph{immutavel}, --- é \emph{infinito}, porque, se fosse \emph{finito}, estaria em potencia para uma perfeição ulterior, e deixaria de ser \emph{imnmutavel}. --- é \emph{simplicissimo}, porque, se fosse \emph{composto}, tambem deixaria de ser \emph{imnmitavel}. --- O ente \emph{necessario}, \emph{infinito}, \emph{simplicissimo}, é \emph{immutavel}; porque, --- sc o ente \emph{necessario} fosse \emph{mwiabcel}, qu adquiriria uma perfeição que não tinha, ou perdéria uma perfeição que tinha, e assim não seria \emph{necessario}, mas \emph{contingente} em relação á perfeição adquirida ou perdida, --- se o ente \emph{infinito} fosse \emph{mutavel}, tambem ou perderia uma perfeição, e não seria \emph{infinito}, mas \emph{finito}, --- s\& o ente \emph{simplicissimo} "fosse mutavel, passaria da potencia para o acto, e já não seria \emph{simplicissimo}, mas \emph{composto} (Cf. \emph{Sum}. \emph{Th}., l. c.). Como se vê, \emph{acto puro}, \emph{ente} necessario, \emph{ente infinito}, \emph{ente simplicissimo}, \emph{ente} immutavel, são attributos que exprimem uma e a mesma coisa, e por isso \& um ente, ao qual convenha um d'esses attributos, convêm todos, e de um d' elles pode legitimamente deduzir-se outro.}.. \textasciicircum{} 88 O ente mutavel é contingente, finito, composto; o ente contingente, finito, composto, é mutavel.

a) \emph{O ente mutavel é contingente}, \emph{finito}, \emph{composto}. --- Ente mutavel é o que está sujeito a mudanças, e por isso não é acto puro. Ora o que não é acto puro é contingente, finito, composto. Logo o ente \emph{mutavel é contingente}, . \emph{finito}, \emph{composto}.

. b) \emph{O ente contingente}, \emph{finito}, \emph{composto}, \emph{é mutavel}. --- Todo o ente contingente, finito, composto, não é acto puro, mas tem em si alguma potencialidade. Ora o que tem potencialidade é mutavel. Logo o ente \emph{contingente}, \emph{finito}, \emph{composto}, é mutavel \footnote{O ente \emph{mutavel} --- é \emph{contingente}, porque, na mudança, adquire ou perde uma pefeic\&o, e assim é \emph{contingente} em relação á perfeição adquirida Ou \emph{perdida}, --- é \emph{finito}, porque, se, na mudarça, adquire ou perde uma perfeição, não possue todas as perfeições, e por isso é \emph{finito}, --- é \emph{composto}, porque, quando \emph{muda}, passa da potencia para o acto, e por isso consta d' esses dois elementos. --- O ente. \emph{contingente}, \emph{finito}, \emph{composto}. \emph{é mutavel}, E mutavel \emph{o ente contingente}; porque, sendo \emph{contingente}. passou e pode passar da potencia para o acto, e assim mudar-se. É nudavel o ente \emph{finito}; porque, sendo \emph{finito}, \emph{node adquirir} uma perfeição ulterior, e assim passar da potencia \emph{para} o \emph{acto}, \emph{É mutavel o} ente \emph{composto}; porque, sendo composto de \emph{potencia} e de \emph{acto}, \emph{pode} passar de um estado para outro. (\emph{Sum}. \emph{Te}, \textbar{}. c., a. 2).}.

89. O ente mutavel suppõe o \textbf{ente immutavel}. --- Na ordem \emph{ontologica}; ou da \emph{actuatidude}, é ente actual deve preceder o ente \emph{possivel} ; porque este só pode adquirir 0 acto pela acção de um 